
\documentclass{article}
\usepackage{colortbl}
\usepackage{makecell}
\usepackage{multirow}
\usepackage{supertabular}

\begin{document}

\newcounter{utterance}

\centering \large Interaction Transcript for game `matchit\_ascii', experiment `same\_grid', episode 4 with reflexion{-}llama{-}3.1{-}8b{-}instruct.
\vspace{24pt}

{ \footnotesize  \setcounter{utterance}{1}
\setlength{\tabcolsep}{0pt}
\begin{supertabular}{c@{$\;$}|p{.15\linewidth}@{}p{.15\linewidth}p{.15\linewidth}p{.15\linewidth}p{.15\linewidth}p{.15\linewidth}}
    \# & $\;$A & \multicolumn{4}{c}{Game Master} & $\;\:$B\\
    \hline

    \theutterance \stepcounter{utterance}  
    & & \multicolumn{4}{p{0.6\linewidth}}{
        \cellcolor[rgb]{0.9,0.9,0.9}{
            \makecell[{{p{\linewidth}}}]{
                \texttt{\tiny{[P1$\langle$GM]}}
                \texttt{You have attempted this same task before and did not succeed. These are the plans you wrote after those attempts.} \\
\texttt{Trial \#1: New plan: } \\
\\ 
\texttt{In my last attempt, I failed to identify the differences between the two grids, which led to a wrong conclusion that they are different grids. Upon reflection, I realized that I didn't ask enough questions to gather more information about the other grid. Specifically, I didn't ask about the number of X's in the middle row, the alignment of the X's in the middle row, or the pattern of ▢'s in the middle row.} \\
\\ 
\texttt{Next time, I will ask more targeted questions to gather more information about the other grid. For example, I will ask about the number of X's in the middle row, the alignment of the X's in the middle row, and the pattern of ▢'s in the middle row. I will also pay closer attention to the answers provided by the Game Master and use that information to inform my next questions.} \\
\\ 
\texttt{Trial \#2: New plan: } \\
\\ 
\texttt{In my last attempt, I failed to identify the differences between the two grids, which led to a wrong conclusion that they are different grids. Upon reflection, I realized that I didn't ask enough questions to gather more information about the other grid. Specifically, I didn't ask about the number of X's in the middle row, the alignment of the X's in the middle row, or the pattern of ▢'s in the middle row.} \\
\\ 
\texttt{However, I also noticed that I didn't use the information provided by the Game Master effectively. For example, when the Game Master said that the X's in the middle row of my grid are not aligned directly above the X's in the middle row of the other grid, I should have asked more follow{-}up questions to clarify the exact nature of the difference.} \\
\\ 
\texttt{Next time, I will ask more targeted questions to gather more information about the other grid, and I will also pay closer attention to the answers provided by the Game Master and use that information to inform my next questions. I will also make sure to ask follow{-}up questions to clarify any ambiguities or inconsistencies in the information provided.} \\
\\ 
\texttt{Specifically, I will ask questions like: "What is the exact number of X's in the middle row of your grid?", "Are the X's in the middle row of your grid aligned directly above the X's in the middle row of my grid, but shifted by one position?", "Are the ▢'s in the middle row of your grid distributed in a specific pattern, such as a diagonal or a horizontal line?".} \\
\\ 
\texttt{By asking more targeted and follow{-}up questions, I hope to gather more information about the other grid and make a more informed decision about whether the two grids are the same or different.} \\
\\ 
\texttt{Trial \#3: New plan: } \\
\\ 
\texttt{In my last attempt, I failed to identify the differences between the two grids, which led to a wrong conclusion that they are different grids. Upon reflection, I realized that I didn't ask enough questions to gather more information about the other grid. Specifically, I didn't ask about the number of X's in the middle row, the alignment of the X's in the middle row, or the pattern of ▢'s in the middle row.} \\
\\ 
\texttt{However, I also noticed that I didn't use the information provided by the Game Master effectively. For example, when the Game Master said that the X's in the middle row of my grid are not aligned directly above the X's in the middle row of the other grid, I should have asked more follow{-}up questions to clarify the exact nature of the difference.} \\
\\ 
\texttt{Next time, I will ask more targeted questions to gather more information about the other grid, and I will also pay closer attention to the answers provided by the Game Master and use that information to inform my next questions. I will also make sure to ask follow{-}up questions to clarify any ambiguities or inconsistencies in the information provided.} \\
\\ 
\texttt{Specifically, I will ask questions like: "What is the exact number of X's in the middle row of your grid?", "Are the X's in the middle row of your grid aligned directly above the X's in the middle row of my grid, but shifted by one position?", "Are the ▢'s in the middle row of your grid distributed in a specific pattern, such as a diagonal or a horizontal line?".} \\
\\ 
\texttt{Additionally, I will also pay closer attention to the Game Master's answers and look for any inconsistencies or contradictions between the answers. If I notice any inconsistencies, I will ask follow{-}up questions to clarify the issue.} \\
\\ 
\texttt{By asking more targeted and follow{-}up questions, and by paying closer attention to the Game Master's answers, I hope to gather more information about the other grid and make a more informed decision about whether the two grids are the same or different.} \\
\\ 
\\ 
\texttt{You are participating in a collaborative guessing game.} \\
\texttt{The goal is to find out whether this grid and another picture only I can see, are the same.} \\
\texttt{Please describe your grid first. Then, I will provide my description and we can ask each other questions about the grids to figure out whether they are the same. Now start your short grid description with "DESCRIPTION:" followed by the description. Do not add anything else. This is your grid:} \\
\texttt{X X X X X} \\
\texttt{▢ ▢ X ▢ ▢} \\
\texttt{▢ ▢ X ▢ ▢} \\
\texttt{▢ ▢ X ▢ ▢} \\
\texttt{▢ ▢ X ▢ ▢} \\
            }
        }
    }
    & \\ \\

    \theutterance \stepcounter{utterance}  
    & \multicolumn{4}{p{0.6\linewidth}}{
        \cellcolor[rgb]{0.9,0.9,0.9}{
            \makecell[{{p{\linewidth}}}]{
                \texttt{\tiny{[P1$\rangle$GM]}}
                \texttt{DESCRIPTION: } \\
\texttt{X X X X X} \\
\texttt{▢ ▢ X ▢ ▢} \\
\texttt{▢ ▢ X ▢ ▢} \\
\texttt{▢ ▢ X ▢ ▢} \\
\texttt{▢ ▢ X ▢ ▢} \\
            }
        }
    }
    & & \\ \\

    \theutterance \stepcounter{utterance}  
    & & & \multicolumn{2}{p{0.3\linewidth}}{
        \cellcolor[rgb]{0.9,0.9,0.9}{
            \makecell[{{p{\linewidth}}}]{
                \texttt{\tiny{[GM$|$GM]}}
                \texttt{continue} \\
            }
        }
    }
    & & \\ \\

    \theutterance \stepcounter{utterance}  
    & & \multicolumn{4}{p{0.6\linewidth}}{
        \cellcolor[rgb]{0.9,0.9,0.9}{
            \makecell[{{p{\linewidth}}}]{
                \texttt{\tiny{[GM$\rangle$P2]}}
                \texttt{You have attempted this same task before and did not succeed. These are the plans you wrote after those attempts.} \\
\texttt{Trial \#1: New plan: } \\
\\ 
\texttt{I got stuck in a loop of asking questions that didn't provide enough new information to make a decision. I started with a good description of my grid, but then my questions were too focused on the ▢'s and didn't explore the X's enough. I also didn't use the information I got from the Game Master's answers effectively. Next time, I will ask more targeted questions that focus on the differences between the grids, and I will use the information I get from the Game Master's answers to inform my next questions. For example, I could ask more questions about the X's, such as their alignment and spacing, and I could ask more questions about the ▢'s, such as their distribution and pattern. I will also try to make a decision earlier, based on the information I have, rather than waiting until the end.} \\
\\ 
\texttt{Trial \#2: New plan: } \\
\\ 
\texttt{I got stuck in a loop of asking questions that didn't provide enough new information to make a decision, and I didn't use the information I got from the Game Master's answers effectively. However, I also made a decision at the end that was based on too little information. Next time, I will ask more targeted questions that focus on the differences between the grids, and I will use the information I get from the Game Master's answers to inform my next questions. I will also try to make a decision earlier, based on the information I have, rather than waiting until the end. } \\
\\ 
\texttt{To achieve this, I will pay closer attention to the Game Master's answers and use them to guide my next questions. For example, when the Game Master said that the ▢'s in their grid are scattered randomly, I could have asked more questions about the X's to see if there are any patterns or differences between the grids. I will also try to ask more open{-}ended questions that encourage the Game Master to provide more information, rather than just asking yes or no questions. } \\
\\ 
\texttt{Additionally, I will try to make a decision earlier, based on the information I have, rather than waiting until the end. This might involve making an educated guess based on the information I have, or asking more questions to clarify any doubts I have. By doing this, I hope to avoid getting stuck in a loop of questions and make a more informed decision about whether the grids are the same or different.} \\
\\ 
\texttt{Trial \#3: New plan: } \\
\\ 
\texttt{I got stuck in a loop of asking questions that didn't provide enough new information to make a decision, and I didn't use the information I got from the Game Master's answers effectively. However, I also made a decision at the end that was based on too little information. Next time, I will ask more targeted questions that focus on the differences between the grids, and I will use the information I get from the Game Master's answers to inform my next questions. I will also try to make a decision earlier, based on the information I have, rather than waiting until the end. } \\
\\ 
\texttt{To achieve this, I will pay closer attention to the Game Master's answers and use them to guide my next questions. For example, when the Game Master said that the top row of their grid consists entirely of ▢'s, I could have asked more questions about the X's in the third column to see if there are any patterns or differences between the grids. I will also try to ask more open{-}ended questions that encourage the Game Master to provide more information, rather than just asking yes or no questions. } \\
\\ 
\texttt{Additionally, I will try to make a decision earlier, based on the information I have, rather than waiting until the end. This might involve making an educated guess based on the information I have, or asking more questions to clarify any doubts I have. By doing this, I hope to avoid getting stuck in a loop of questions and make a more informed decision about whether the grids are the same or different.} \\
\\ 
\\ 
\texttt{You are participating in a collaborative guessing game.} \\
\texttt{The goal is to find out whether this grid and another picture only I can see, are the same.} \\
\texttt{Please describe your grid first. Then, I will provide my description and we can ask each other questions about the grids to figure out whether they are the same. Now start your short grid description with "DESCRIPTION:" followed by the description. Do not add anything else. This is your grid:} \\
\texttt{X X X X X} \\
\texttt{▢ ▢ X ▢ ▢} \\
\texttt{▢ ▢ X ▢ ▢} \\
\texttt{▢ ▢ X ▢ ▢} \\
\texttt{▢ ▢ X ▢ ▢} \\
            }
        }
    }
    & \\ \\

    \theutterance \stepcounter{utterance}  
    & & & \multicolumn{4}{p{0.6\linewidth}}{
        \cellcolor[rgb]{0.9,0.9,0.9}{
            \makecell[{{p{\linewidth}}}]{
                \texttt{\tiny{[GM$\langle$P2]}}
                \texttt{DESCRIPTION: My grid is a 5x5 grid with a pattern of X's and ▢'s. The top row consists entirely of X's. The second, third, and fourth rows each have two ▢'s and three X's, with the X's in the third column. The bottom row consists entirely of ▢'s.} \\
            }
        }
    }
     \\ \\

    \theutterance \stepcounter{utterance}  
    & & & \multicolumn{2}{p{0.3\linewidth}}{
        \cellcolor[rgb]{0.9,0.9,0.9}{
            \makecell[{{p{\linewidth}}}]{
                \texttt{\tiny{[GM$|$GM]}}
                \texttt{continue} \\
            }
        }
    }
    & & \\ \\

    \theutterance \stepcounter{utterance}  
    & & \multicolumn{4}{p{0.6\linewidth}}{
        \cellcolor[rgb]{0.9,0.9,0.9}{
            \makecell[{{p{\linewidth}}}]{
                \texttt{\tiny{[GM$\rangle$P2]}}
                \texttt{This is my DESCRIPTION: } \\
\texttt{X X X X X} \\
\texttt{▢ ▢ X ▢ ▢} \\
\texttt{▢ ▢ X ▢ ▢} \\
\texttt{▢ ▢ X ▢ ▢} \\
\texttt{▢ ▢ X ▢ ▢} \\
\texttt{Now ask a question in order to find out new aspects of my grid that may be different to your grid. Start with "QUESTION:" and do not add anything else.} \\
            }
        }
    }
    & \\ \\

    \theutterance \stepcounter{utterance}  
    & & & \multicolumn{4}{p{0.6\linewidth}}{
        \cellcolor[rgb]{0.9,0.9,0.9}{
            \makecell[{{p{\linewidth}}}]{
                \texttt{\tiny{[GM$\langle$P2]}}
                \texttt{QUESTION: Are there any ▢'s in the first column of your grid?} \\
            }
        }
    }
     \\ \\

    \theutterance \stepcounter{utterance}  
    & & & \multicolumn{2}{p{0.3\linewidth}}{
        \cellcolor[rgb]{0.9,0.9,0.9}{
            \makecell[{{p{\linewidth}}}]{
                \texttt{\tiny{[GM$|$GM]}}
                \texttt{continue} \\
            }
        }
    }
    & & \\ \\

    \theutterance \stepcounter{utterance}  
    & & \multicolumn{4}{p{0.6\linewidth}}{
        \cellcolor[rgb]{0.9,0.9,0.9}{
            \makecell[{{p{\linewidth}}}]{
                \texttt{\tiny{[P1$\langle$GM]}}
                \texttt{This is my DESCRIPTION: My grid is a 5x5 grid with a pattern of X's and ▢'s. The top row consists entirely of X's. The second, third, and fourth rows each have two ▢'s and three X's, with the X's in the third column. The bottom row consists entirely of ▢'s.} \\
\texttt{QUESTION: Are there any ▢'s in the first column of your grid? Start your answer with "ANSWER:" and do not add anything else.} \\
            }
        }
    }
    & \\ \\

    \theutterance \stepcounter{utterance}  
    & \multicolumn{4}{p{0.6\linewidth}}{
        \cellcolor[rgb]{0.9,0.9,0.9}{
            \makecell[{{p{\linewidth}}}]{
                \texttt{\tiny{[P1$\rangle$GM]}}
                \texttt{ANSWER: No, there are no ▢'s in the first column of my grid.} \\
            }
        }
    }
    & & \\ \\

    \theutterance \stepcounter{utterance}  
    & & & \multicolumn{2}{p{0.3\linewidth}}{
        \cellcolor[rgb]{0.9,0.9,0.9}{
            \makecell[{{p{\linewidth}}}]{
                \texttt{\tiny{[GM$|$GM]}}
                \texttt{continue} \\
            }
        }
    }
    & & \\ \\

    \theutterance \stepcounter{utterance}  
    & & \multicolumn{4}{p{0.6\linewidth}}{
        \cellcolor[rgb]{0.9,0.9,0.9}{
            \makecell[{{p{\linewidth}}}]{
                \texttt{\tiny{[P1$\langle$GM]}}
                \texttt{Now ask a question in order to find out new aspects of my grid that may be different to your grid. Start with "QUESTION:" and do not add anything else.} \\
            }
        }
    }
    & \\ \\

    \theutterance \stepcounter{utterance}  
    & \multicolumn{4}{p{0.6\linewidth}}{
        \cellcolor[rgb]{0.9,0.9,0.9}{
            \makecell[{{p{\linewidth}}}]{
                \texttt{\tiny{[P1$\rangle$GM]}}
                \texttt{QUESTION: Are the X's in the third column of your grid aligned directly above the X's in the third column of my grid?} \\
            }
        }
    }
    & & \\ \\

    \theutterance \stepcounter{utterance}  
    & & & \multicolumn{2}{p{0.3\linewidth}}{
        \cellcolor[rgb]{0.9,0.9,0.9}{
            \makecell[{{p{\linewidth}}}]{
                \texttt{\tiny{[GM$|$GM]}}
                \texttt{continue} \\
            }
        }
    }
    & & \\ \\

    \theutterance \stepcounter{utterance}  
    & & \multicolumn{4}{p{0.6\linewidth}}{
        \cellcolor[rgb]{0.9,0.9,0.9}{
            \makecell[{{p{\linewidth}}}]{
                \texttt{\tiny{[GM$\rangle$P2]}}
                \texttt{ANSWER: No, there are no ▢'s in the first column of my grid.} \\
\texttt{QUESTION: Are the X's in the third column of your grid aligned directly above the X's in the third column of my grid? Start your answer with "ANSWER:" and do not add anything else.} \\
            }
        }
    }
    & \\ \\

    \theutterance \stepcounter{utterance}  
    & & & \multicolumn{4}{p{0.6\linewidth}}{
        \cellcolor[rgb]{0.9,0.9,0.9}{
            \makecell[{{p{\linewidth}}}]{
                \texttt{\tiny{[GM$\langle$P2]}}
                \texttt{ANSWER: Yes, the X's in the third column of your grid are aligned directly above the X's in the third column of my grid.} \\
            }
        }
    }
     \\ \\

    \theutterance \stepcounter{utterance}  
    & & & \multicolumn{2}{p{0.3\linewidth}}{
        \cellcolor[rgb]{0.9,0.9,0.9}{
            \makecell[{{p{\linewidth}}}]{
                \texttt{\tiny{[GM$|$GM]}}
                \texttt{continue} \\
            }
        }
    }
    & & \\ \\

    \theutterance \stepcounter{utterance}  
    & & \multicolumn{4}{p{0.6\linewidth}}{
        \cellcolor[rgb]{0.9,0.9,0.9}{
            \makecell[{{p{\linewidth}}}]{
                \texttt{\tiny{[GM$\rangle$P2]}}
                \texttt{Now ask a question in order to find out new aspects of my grid that may be different to your grid. Start with "QUESTION:" and do not add anything else.} \\
            }
        }
    }
    & \\ \\

    \theutterance \stepcounter{utterance}  
    & & & \multicolumn{4}{p{0.6\linewidth}}{
        \cellcolor[rgb]{0.9,0.9,0.9}{
            \makecell[{{p{\linewidth}}}]{
                \texttt{\tiny{[GM$\langle$P2]}}
                \texttt{QUESTION: Are there any rows in your grid that consist entirely of ▢'s?} \\
            }
        }
    }
     \\ \\

    \theutterance \stepcounter{utterance}  
    & & & \multicolumn{2}{p{0.3\linewidth}}{
        \cellcolor[rgb]{0.9,0.9,0.9}{
            \makecell[{{p{\linewidth}}}]{
                \texttt{\tiny{[GM$|$GM]}}
                \texttt{continue} \\
            }
        }
    }
    & & \\ \\

    \theutterance \stepcounter{utterance}  
    & & \multicolumn{4}{p{0.6\linewidth}}{
        \cellcolor[rgb]{0.9,0.9,0.9}{
            \makecell[{{p{\linewidth}}}]{
                \texttt{\tiny{[P1$\langle$GM]}}
                \texttt{ANSWER: Yes, the X's in the third column of your grid are aligned directly above the X's in the third column of my grid.} \\
\texttt{QUESTION: Are there any rows in your grid that consist entirely of ▢'s? Start your answer with "ANSWER:" and do not add anything else.} \\
            }
        }
    }
    & \\ \\

    \theutterance \stepcounter{utterance}  
    & \multicolumn{4}{p{0.6\linewidth}}{
        \cellcolor[rgb]{0.9,0.9,0.9}{
            \makecell[{{p{\linewidth}}}]{
                \texttt{\tiny{[P1$\rangle$GM]}}
                \texttt{ANSWER: Yes, the bottom row of my grid consists entirely of ▢'s.} \\
            }
        }
    }
    & & \\ \\

    \theutterance \stepcounter{utterance}  
    & & & \multicolumn{2}{p{0.3\linewidth}}{
        \cellcolor[rgb]{0.9,0.9,0.9}{
            \makecell[{{p{\linewidth}}}]{
                \texttt{\tiny{[GM$|$GM]}}
                \texttt{continue} \\
            }
        }
    }
    & & \\ \\

    \theutterance \stepcounter{utterance}  
    & & \multicolumn{4}{p{0.6\linewidth}}{
        \cellcolor[rgb]{0.9,0.9,0.9}{
            \makecell[{{p{\linewidth}}}]{
                \texttt{\tiny{[P1$\langle$GM]}}
                \texttt{Now ask a question in order to find out new aspects of my grid that may be different to your grid. Start with "QUESTION:" and do not add anything else.} \\
            }
        }
    }
    & \\ \\

    \theutterance \stepcounter{utterance}  
    & \multicolumn{4}{p{0.6\linewidth}}{
        \cellcolor[rgb]{0.9,0.9,0.9}{
            \makecell[{{p{\linewidth}}}]{
                \texttt{\tiny{[P1$\rangle$GM]}}
                \texttt{QUESTION: Are there any rows in your grid that consist entirely of X's?} \\
            }
        }
    }
    & & \\ \\

    \theutterance \stepcounter{utterance}  
    & & & \multicolumn{2}{p{0.3\linewidth}}{
        \cellcolor[rgb]{0.9,0.9,0.9}{
            \makecell[{{p{\linewidth}}}]{
                \texttt{\tiny{[GM$|$GM]}}
                \texttt{continue} \\
            }
        }
    }
    & & \\ \\

    \theutterance \stepcounter{utterance}  
    & & \multicolumn{4}{p{0.6\linewidth}}{
        \cellcolor[rgb]{0.9,0.9,0.9}{
            \makecell[{{p{\linewidth}}}]{
                \texttt{\tiny{[GM$\rangle$P2]}}
                \texttt{ANSWER: Yes, the bottom row of my grid consists entirely of ▢'s.} \\
\texttt{QUESTION: Are there any rows in your grid that consist entirely of X's? Start your answer with "ANSWER:" and do not add anything else.} \\
            }
        }
    }
    & \\ \\

    \theutterance \stepcounter{utterance}  
    & & & \multicolumn{4}{p{0.6\linewidth}}{
        \cellcolor[rgb]{0.9,0.9,0.9}{
            \makecell[{{p{\linewidth}}}]{
                \texttt{\tiny{[GM$\langle$P2]}}
                \texttt{ANSWER: No, there are no rows in my grid that consist entirely of X's.} \\
            }
        }
    }
     \\ \\

    \theutterance \stepcounter{utterance}  
    & & & \multicolumn{2}{p{0.3\linewidth}}{
        \cellcolor[rgb]{0.9,0.9,0.9}{
            \makecell[{{p{\linewidth}}}]{
                \texttt{\tiny{[GM$|$GM]}}
                \texttt{continue} \\
            }
        }
    }
    & & \\ \\

    \theutterance \stepcounter{utterance}  
    & & \multicolumn{4}{p{0.6\linewidth}}{
        \cellcolor[rgb]{0.9,0.9,0.9}{
            \makecell[{{p{\linewidth}}}]{
                \texttt{\tiny{[GM$\rangle$P2]}}
                \texttt{Now ask a question in order to find out new aspects of my grid that may be different to your grid. Start with "QUESTION:" and do not add anything else.} \\
            }
        }
    }
    & \\ \\

    \theutterance \stepcounter{utterance}  
    & & & \multicolumn{4}{p{0.6\linewidth}}{
        \cellcolor[rgb]{0.9,0.9,0.9}{
            \makecell[{{p{\linewidth}}}]{
                \texttt{\tiny{[GM$\langle$P2]}}
                \texttt{QUESTION: Are there any columns in your grid that consist entirely of X's?} \\
            }
        }
    }
     \\ \\

    \theutterance \stepcounter{utterance}  
    & & & \multicolumn{2}{p{0.3\linewidth}}{
        \cellcolor[rgb]{0.9,0.9,0.9}{
            \makecell[{{p{\linewidth}}}]{
                \texttt{\tiny{[GM$|$GM]}}
                \texttt{continue} \\
            }
        }
    }
    & & \\ \\

    \theutterance \stepcounter{utterance}  
    & & \multicolumn{4}{p{0.6\linewidth}}{
        \cellcolor[rgb]{0.9,0.9,0.9}{
            \makecell[{{p{\linewidth}}}]{
                \texttt{\tiny{[P1$\langle$GM]}}
                \texttt{ANSWER: No, there are no rows in my grid that consist entirely of X's.} \\
\texttt{QUESTION: Are there any columns in your grid that consist entirely of X's? Start your answer with "ANSWER:" and do not add anything else.} \\
            }
        }
    }
    & \\ \\

    \theutterance \stepcounter{utterance}  
    & \multicolumn{4}{p{0.6\linewidth}}{
        \cellcolor[rgb]{0.9,0.9,0.9}{
            \makecell[{{p{\linewidth}}}]{
                \texttt{\tiny{[P1$\rangle$GM]}}
                \texttt{ANSWER: No, there are no columns in my grid that consist entirely of X's.} \\
            }
        }
    }
    & & \\ \\

    \theutterance \stepcounter{utterance}  
    & & & \multicolumn{2}{p{0.3\linewidth}}{
        \cellcolor[rgb]{0.9,0.9,0.9}{
            \makecell[{{p{\linewidth}}}]{
                \texttt{\tiny{[GM$|$GM]}}
                \texttt{continue} \\
            }
        }
    }
    & & \\ \\

    \theutterance \stepcounter{utterance}  
    & & \multicolumn{4}{p{0.6\linewidth}}{
        \cellcolor[rgb]{0.9,0.9,0.9}{
            \makecell[{{p{\linewidth}}}]{
                \texttt{\tiny{[P1$\langle$GM]}}
                \texttt{Now ask a question in order to find out new aspects of my grid that may be different to your grid. Start with "QUESTION:" and do not add anything else.} \\
            }
        }
    }
    & \\ \\

    \theutterance \stepcounter{utterance}  
    & \multicolumn{4}{p{0.6\linewidth}}{
        \cellcolor[rgb]{0.9,0.9,0.9}{
            \makecell[{{p{\linewidth}}}]{
                \texttt{\tiny{[P1$\rangle$GM]}}
                \texttt{QUESTION: Are there any rows in your grid that have a different number of X's and ▢'s than the rows in my grid?} \\
            }
        }
    }
    & & \\ \\

    \theutterance \stepcounter{utterance}  
    & & & \multicolumn{2}{p{0.3\linewidth}}{
        \cellcolor[rgb]{0.9,0.9,0.9}{
            \makecell[{{p{\linewidth}}}]{
                \texttt{\tiny{[GM$|$GM]}}
                \texttt{continue} \\
            }
        }
    }
    & & \\ \\

    \theutterance \stepcounter{utterance}  
    & & \multicolumn{4}{p{0.6\linewidth}}{
        \cellcolor[rgb]{0.9,0.9,0.9}{
            \makecell[{{p{\linewidth}}}]{
                \texttt{\tiny{[GM$\rangle$P2]}}
                \texttt{ANSWER: No, there are no columns in my grid that consist entirely of X's.} \\
\texttt{QUESTION: Are there any rows in your grid that have a different number of X's and ▢'s than the rows in my grid? Start your answer with "ANSWER:" and do not add anything else.} \\
            }
        }
    }
    & \\ \\

    \theutterance \stepcounter{utterance}  
    & & & \multicolumn{4}{p{0.6\linewidth}}{
        \cellcolor[rgb]{0.9,0.9,0.9}{
            \makecell[{{p{\linewidth}}}]{
                \texttt{\tiny{[GM$\langle$P2]}}
                \texttt{ANSWER: Yes, the top row of my grid has 5 X's and 0 ▢'s, but the top row of your grid has 5 X's and 0 ▢'s as well. However, the bottom row of your grid has 0 X's and 5 ▢'s, which is different from the bottom row of my grid, which has 5 ▢'s and 0 X's.} \\
            }
        }
    }
     \\ \\

    \theutterance \stepcounter{utterance}  
    & & & \multicolumn{2}{p{0.3\linewidth}}{
        \cellcolor[rgb]{0.9,0.9,0.9}{
            \makecell[{{p{\linewidth}}}]{
                \texttt{\tiny{[GM$|$GM]}}
                \texttt{continue} \\
            }
        }
    }
    & & \\ \\

    \theutterance \stepcounter{utterance}  
    & & \multicolumn{4}{p{0.6\linewidth}}{
        \cellcolor[rgb]{0.9,0.9,0.9}{
            \makecell[{{p{\linewidth}}}]{
                \texttt{\tiny{[GM$\rangle$P2]}}
                \texttt{Now come to a decision. What do you think: are your grid and the other grid described the same?  Write "DECISION: same grid" if you think they are the same grid or "DECISION: different grids" if you think they are different grids. Do not add anything else.} \\
            }
        }
    }
    & \\ \\

    \theutterance \stepcounter{utterance}  
    & & & \multicolumn{4}{p{0.6\linewidth}}{
        \cellcolor[rgb]{0.9,0.9,0.9}{
            \makecell[{{p{\linewidth}}}]{
                \texttt{\tiny{[GM$\langle$P2]}}
                \texttt{DECISION: different grids} \\
            }
        }
    }
     \\ \\

    \theutterance \stepcounter{utterance}  
    & & & \multicolumn{2}{p{0.3\linewidth}}{
        \cellcolor[rgb]{0.9,0.9,0.9}{
            \makecell[{{p{\linewidth}}}]{
                \texttt{\tiny{[GM$|$GM]}}
                \texttt{continue} \\
            }
        }
    }
    & & \\ \\

    \theutterance \stepcounter{utterance}  
    & & & \multicolumn{2}{p{0.3\linewidth}}{
        \cellcolor[rgb]{0.9,0.9,0.9}{
            \makecell[{{p{\linewidth}}}]{
                \texttt{\tiny{[GM$|$GM]}}
                \texttt{loss} \\
            }
        }
    }
    & & \\ \\

    \theutterance \stepcounter{utterance}  
    & & \multicolumn{4}{p{0.6\linewidth}}{
        \cellcolor[rgb]{0.9,0.9,0.9}{
            \makecell[{{p{\linewidth}}}]{
                \texttt{\tiny{[P1$\langle$GM]}}
                \texttt{ANSWER: Yes, the top row of my grid has 5 X's and 0 ▢'s, but the top row of your grid has 5 X's and 0 ▢'s as well. However, the bottom row of your grid has 0 X's and 5 ▢'s, which is different from the bottom row of my grid, which has 5 ▢'s and 0 X's.} \\
\texttt{Now come to a decision. What do you think: are your grid and the other grid described the same?  Write "DECISION: same grid" if you think they are the same grid or "DECISION: different grids" if you think they are different grids. Do not add anything else.} \\
            }
        }
    }
    & \\ \\

    \theutterance \stepcounter{utterance}  
    & \multicolumn{4}{p{0.6\linewidth}}{
        \cellcolor[rgb]{0.9,0.9,0.9}{
            \makecell[{{p{\linewidth}}}]{
                \texttt{\tiny{[P1$\rangle$GM]}}
                \texttt{DECISION: different grids} \\
            }
        }
    }
    & & \\ \\

    \theutterance \stepcounter{utterance}  
    & & & \multicolumn{2}{p{0.3\linewidth}}{
        \cellcolor[rgb]{0.9,0.9,0.9}{
            \makecell[{{p{\linewidth}}}]{
                \texttt{\tiny{[GM$|$GM]}}
                \texttt{continue} \\
            }
        }
    }
    & & \\ \\

    \theutterance \stepcounter{utterance}  
    & & & \multicolumn{2}{p{0.3\linewidth}}{
        \cellcolor[rgb]{0.9,0.9,0.9}{
            \makecell[{{p{\linewidth}}}]{
                \texttt{\tiny{[GM$|$GM]}}
                \texttt{loss} \\
            }
        }
    }
    & & \\ \\

    \theutterance \stepcounter{utterance}  
    & \multicolumn{4}{p{0.6\linewidth}}{
        \cellcolor[rgb]{0.95,0.95,0.95}{
            \makecell[{{p{\linewidth}}}]{
                \texttt{\tiny{[P2$\rightarrow$P1]}}
                \texttt{Post{-}episode Reflexion (not sent to GM):} \\
\texttt{New plan: } \\
\\ 
\texttt{In my last attempts, I failed to identify the differences between the two grids, which led to a wrong conclusion that they are different grids. Upon reflection, I realized that I didn't ask enough questions to gather more information about the other grid. Specifically, I didn't ask about the number of X's in the middle row, the alignment of the X's in the middle row, or the pattern of ▢'s in the middle row.} \\
\\ 
\texttt{However, I also noticed that I didn't use the information provided by the Game Master effectively. For example, when the Game Master said that the X's in the middle row of my grid are not aligned directly above the X's in the middle row of the other grid, I should have asked more follow{-}up questions to clarify the exact nature of the difference.} \\
\\ 
\texttt{Next time, I will ask more targeted questions to gather more information about the other grid, and I will also pay closer attention to the answers provided by the Game Master and use that information to inform my next questions. I will also make sure to ask follow{-}up questions to clarify any ambiguities or inconsistencies in the information provided.} \\
\\ 
\texttt{Specifically, I will ask questions like: "What is the exact number of X's in the middle row of your grid?", "Are the X's in the middle row of your grid aligned directly above the X's in the middle row of my grid, but shifted by one position?", "Are the ▢'s in the middle row of your grid distributed in a specific pattern, such as a diagonal or a horizontal line?".} \\
\\ 
\texttt{Additionally, I will also pay closer attention to the Game Master's answers and look for any inconsistencies or contradictions between the answers. If I notice any inconsistencies, I will ask follow{-}up questions to clarify the issue.} \\
\\ 
\texttt{However, I also noticed that I didn't ask about the top row of the other grid, which is entirely X's. I should have asked if the top row of the other grid is also entirely X's. This would have given me more information about the other grid.} \\
\\ 
\texttt{Next time, I will ask about the top row of the other grid, and I will also ask about the bottom row of the other grid, which is entirely ▢'s. I will ask questions like: "Is the top row of your grid also entirely X's?", "Is the bottom row of your grid also entirely ▢'s?".} \\
\\ 
\texttt{By asking more targeted and follow{-}up questions, and by paying closer attention to the Game Master's answers, I hope to gather more information about the other grid and make a more informed decision about whether the two grids are the same or different.} \\
            }
        }
    }
    & & \\ \\

    \theutterance \stepcounter{utterance}  
    & & & \multicolumn{4}{p{0.6\linewidth}}{
        \cellcolor[rgb]{0.95,0.95,0.95}{
            \makecell[{{p{\linewidth}}}]{
                \texttt{\tiny{[P1$\rightarrow$P2]}}
                \texttt{Post{-}episode Reflexion (not sent to GM):} \\
\texttt{New plan: } \\
\\ 
\texttt{I got stuck in a loop of asking questions that didn't provide enough new information to make a decision, and I didn't use the information I got from the Game Master's answers effectively. However, I also made a decision at the end that was based on too little information. Next time, I will ask more targeted questions that focus on the differences between the grids, and I will use the information I get from the Game Master's answers to inform my next questions. I will also try to make a decision earlier, based on the information I have, rather than waiting until the end. } \\
\\ 
\texttt{To achieve this, I will pay closer attention to the Game Master's answers and use them to guide my next questions. For example, when the Game Master said that the ▢'s in their grid are scattered randomly, I could have asked more questions about the X's to see if there are any patterns or differences between the grids. I will also try to ask more open{-}ended questions that encourage the Game Master to provide more information, rather than just asking yes or no questions. } \\
\\ 
\texttt{Additionally, I will try to make a decision earlier, based on the information I have, rather than waiting until the end. This might involve making an educated guess based on the information I have, or asking more questions to clarify any doubts I have. By doing this, I hope to avoid getting stuck in a loop of questions and make a more informed decision about whether the grids are the same or different.} \\
\\ 
\texttt{However, I also noticed that I didn't use the information I got from the Game Master's answers effectively. For example, when the Game Master said that the bottom row of their grid consists entirely of ▢'s, I could have used this information to ask more targeted questions about the X's in the other rows. I will try to use the information I get from the Game Master's answers more effectively in my next questions.} \\
            }
        }
    }
     \\ \\

\end{supertabular}
}

\end{document}
